% CSE 719 — Distributed & Cloud Computing
% Lecture 3: MapReduce & Hadoop — Model Answers (2021–2024)
% University of Chittagong, 8th Semester B.Sc. Engineering
% Note: 2020 paper has no MapReduce questions.
% Compile: pdflatex MapReduce_Solutions.tex (twice for TOC)

\documentclass[12pt, a4paper]{article}

\usepackage[left=2.5cm, right=2.5cm, top=2.8cm, bottom=2.8cm, headheight=15pt]{geometry}
\usepackage{amsmath, amssymb}
\usepackage{booktabs}
\usepackage{array}
\usepackage{xcolor}
\usepackage{enumitem}
\usepackage{fancyhdr}
\usepackage{titlesec}
\usepackage{mdframed}
\usepackage{tabularx}
\usepackage{hyperref}

%----- Colours -----
\definecolor{sectionblue}{RGB}{10,60,110}
\definecolor{boxbg}{RGB}{255,248,240}
\definecolor{answerbg}{RGB}{240,247,255}
\definecolor{notesbg}{RGB}{255,250,230}

%----- Box environments -----
\mdfdefinestyle{qframe}{linecolor=orange!75!black, backgroundcolor=boxbg, roundcorner=4pt, linewidth=1.3pt, innertopmargin=6pt, innerbottommargin=6pt, innerleftmargin=9pt, innerrightmargin=9pt}
\mdfdefinestyle{solframe}{linecolor=sectionblue, backgroundcolor=answerbg, roundcorner=4pt, linewidth=1.3pt, innertopmargin=7pt, innerbottommargin=7pt, innerleftmargin=9pt, innerrightmargin=9pt}
\mdfdefinestyle{notesframe}{linecolor=orange!70!black, backgroundcolor=notesbg, roundcorner=4pt, linewidth=1pt, innertopmargin=5pt, innerbottommargin=5pt, innerleftmargin=8pt, innerrightmargin=8pt}
\newenvironment{qbox}{\begin{mdframed}[style=qframe]}{\end{mdframed}}
\newenvironment{solbox}{\begin{mdframed}[style=solframe]}{\end{mdframed}}
\newenvironment{notesbox}{\begin{mdframed}[style=notesframe]}{\end{mdframed}}

%----- Page style -----
\pagestyle{fancy}
\fancyhf{}
\lhead{\small\color{sectionblue} CSE 719 --- Lecture 3: MapReduce \& Hadoop}
\rhead{\small\color{sectionblue} Model Solutions}
\cfoot{\small\thepage}
\renewcommand{\headrulewidth}{0.5pt}

%----- Section formatting -----
\titleformat{\section}{\Large\bfseries\color{sectionblue}}{}{0em}{}[\vspace{-2pt}\color{sectionblue}\rule{\linewidth}{0.8pt}]
\titleformat{\subsection}{\large\bfseries\color{sectionblue!85!black}}{}{0em}{}

\newcommand{\qmarks}[1]{\hfill\textbf{\color{orange!70!black}[#1 marks]}}

%=====================================================================
\begin{document}
%=====================================================================
\begin{center}
  {\LARGE\bfseries\color{sectionblue} CSE 719: Distributed \& Cloud Computing}\\[0.5em]
  {\Large\bfseries Lecture 3 --- MapReduce \& Hadoop}\\[0.3em]
  {\large Model Answers: 2021, 2022 \& 2024 Papers}\\[0.2em]
  {\small University of Chittagong \quad|\quad 8\textsuperscript{th} Sem B.Sc.\ (Engg.)
  \quad|\quad Source: Lecture-03 slides}
\end{center}
\vspace{0.3em}\noindent\rule{\linewidth}{1pt}\vspace{0.5em}

\noindent\textit{Note: 2020 paper has no MapReduce questions.
In 2024, all four parts of Q1 (Section A) are MapReduce --- treat this as a core topic.}

\tableofcontents
\newpage

%=====================================================================
\section{MapReduce Essentials (Reference)}
%=====================================================================

\begin{solbox}
\textbf{MapReduce:} A programming framework for \textbf{distributed and parallel processing on
large datasets}. Consists of two phases: \textbf{Map} and \textbf{Reduce} (barrier between them:
no Reduce starts until all Maps complete). Terms borrowed from functional programming (Lisp).\\[6pt]

\textbf{Execution pipeline:}
\[
\text{Input (DFS)} \;\xrightarrow{\text{Map}}\; \text{(key,value) pairs (local disk)}
\;\xrightarrow{\text{Shuffle/Sort}}\; \text{Reduce tasks (group by key)}
\;\xrightarrow{\text{Reduce}}\; \text{Output (DFS)}
\]

\textbf{Map:} Called once per input record. Emits zero or more \textbf{intermediate (key, value)}
pairs. Each Map task is \textbf{completely independent} of every other Map task.\\
\textbf{Shuffle:} Framework groups all intermediate values with the same key and routes them to
one Reduce task. \textbf{Hash partitioning:} key $\rightarrow$ reduce\# = hash(key) \% $R$
(number of reducers).\\
\textbf{Reduce:} Receives all values for one key; merges them into output. Each Reduce task
operates independently.\\[6pt]

\textbf{Storage during execution:}
\begin{itemize}[leftmargin=1.4cm, itemsep=1pt]
  \item Map input: \textbf{from distributed FS} (GFS/HDFS).
  \item Map output: \textbf{to local disk} of the Map node.
  \item Reduce input: \textbf{pulled from local disks of Map nodes} (network read).
  \item Reduce output: \textbf{to distributed FS}.
\end{itemize}

\textbf{WordCount (running example):} Input \texttt{<filename, "Welcome Everyone Hello Everyone">}
$\xrightarrow{\text{Map}}$ \texttt{(Welcome,1),(Everyone,1),(Hello,1),(Everyone,1)}
$\xrightarrow{\text{Reduce}}$ \texttt{(Everyone,2),(Hello,1),(Welcome,1)}.
\end{solbox}

%=====================================================================
\section{2021 Examination (CSE 813, 52.5 marks)}
%=====================================================================

%-------
\subsection{Q7(c) --- Various types of MapReduce jobs
  \texorpdfstring{\qmarks{2}}{}}
\begin{qbox}What are the various types of MapReduce jobs?\end{qbox}
\begin{solbox}
The following are standard types of jobs that naturally fit the MapReduce model:

\begin{center}\renewcommand{\arraystretch}{1.3}
\begin{tabularx}{\linewidth}{l X X}
\toprule
\textbf{Job type} & \textbf{Map phase} & \textbf{Reduce phase} \\
\midrule
\textbf{Word count} & Emit \texttt{(word, 1)} for each word & Sum counts per word \\
\textbf{Distributed grep} & Emit line if it matches the pattern & Copy to output (identity) \\
\textbf{Reverse web-link graph} & For each link $a\!\to\!b$, emit \texttt{(b, a)} & Emit \texttt{(b, list(a))} per target page \\
\textbf{URL access frequency} & Emit \texttt{(URL, 1)} per log entry & Sum per URL; chain 2\textsuperscript{nd} MR job for \% \\
\textbf{Distributed sort} & Emit \texttt{(value, \_)} identity & Identity; range partitioning across reducers \\
\textbf{Inverted index} & Emit \texttt{(word, docID)} & Emit \texttt{(word, [docID list])} for search index \\
\bottomrule
\end{tabularx}\end{center}

\textbf{Common characteristics:} large-scale data-intensive, embarrassingly parallel at the record
level, naturally split into independent Map tasks. Particularly suited for \textbf{log processing,
web search indexing}, and \textbf{ad-hoc queries}.
\end{solbox}

%=====================================================================
\section{2022 Examination (CSE 513, 54 marks)}
%=====================================================================

%-------
\subsection{Q4(a) --- Two limitations on the Map function
  \texorpdfstring{\qmarks{2}}{}}
\begin{qbox}What are the two important limitations MapReduce places on the Map function?\end{qbox}
\begin{solbox}
\begin{enumerate}[leftmargin=1.6cm, itemsep=5pt]
  \item \textbf{No side effects / Must be a pure function.}\\
    The Map function must be \textbf{deterministic}: given the same input record, it must always
    produce the same output. It cannot write to any \textbf{shared external state} (databases,
    counters, files) or communicate with other Map tasks during execution.\\
    \textit{Why:} Failed Map tasks are silently \textbf{re-executed} on another worker. If Map had
    side effects, re-execution would corrupt external state (e.g., a database counter would be
    incremented twice). Determinism also enables \textbf{speculative execution} (two copies of the
    same task; use whichever finishes first).

  \item \textbf{No communication between Map tasks / Each record processed independently.}\\
    A Map function sees \textbf{exactly one input record at a time}. It cannot read the outputs of
    other Map tasks, access sibling records, or maintain state across multiple Map invocations.
    Each Map task operates on its own input split in complete \textbf{isolation}.\\
    \textit{Why:} This independence is what makes MapReduce \textbf{trivially parallelisable} ---
    all Map tasks can run simultaneously with no synchronisation overhead.
\end{enumerate}
\end{solbox}

%-------
\subsection{Q4(b) --- How MapReduce handles straggler tasks
  \texorpdfstring{\qmarks{2}}{}}
\begin{qbox}How does MapReduce handle straggler tasks?\end{qbox}
\begin{solbox}
\textbf{Straggler:} A task that runs significantly slower than its peers, delaying the entire job.
Caused by: bad disk, network congestion, CPU throttling, or memory pressure on that node. Because
of the \textbf{barrier} between Map and Reduce phases (no Reduce starts until all Maps finish),
even one slow Map task stalls the whole job.\\[6pt]

\textbf{Solution: Speculative Execution (Backup Tasks)}
\begin{enumerate}[leftmargin=1.6cm, itemsep=3pt]
  \item The framework \textbf{tracks the progress} of each task (percentage complete).
  \item When a job is \textbf{near completion} and some tasks are progressing abnormally slowly,
    the master launches \textbf{backup copies} of those straggler tasks on other available workers.
  \item \textbf{Whichever instance finishes first} (original or backup) has its output accepted;
    the other is killed.
  \item Result: tail latency is greatly reduced at the cost of a small increase in cluster
    resource usage ($\sim$a few percent extra tasks).
\end{enumerate}

\textit{Key insight:} Because Map functions are deterministic and have no side effects, running two
copies of the same task is \textbf{always safe} --- both will produce identical output.
\end{solbox}

%=====================================================================
\section{2024 Examination (CSE-813, 54 marks)}
%=====================================================================

%-------
\subsection{Q1(a) --- HPC fault tolerance using message passing
  \texorpdfstring{\qmarks{2.5}}{}}
\begin{qbox}Describe the fault tolerance technique used in High Performance Computing (HPC)
systems that use message passing.\end{qbox}
\begin{solbox}
Traditional HPC systems (using \textbf{MPI} --- Message Passing Interface) rely on
\textbf{Checkpoint/Restart} for fault tolerance:

\begin{enumerate}[leftmargin=1.6cm, itemsep=3pt]
  \item \textbf{Periodic global checkpointing:} At regular intervals, \textbf{all} MPI processes
    synchronously \textbf{save their entire memory state} (variables, buffers, message queues) to
    stable storage (parallel file system).
  \item \textbf{Restart on failure:} When any process fails, the \textbf{entire application}
    (all processes) rolls back to the last checkpoint and resumes from that saved state.
  \item \textbf{Coordinated recovery:} All processes restart together to ensure a
    \textbf{consistent global state} (no dangling messages, no inconsistency between processes).
\end{enumerate}

\textbf{Limitations at scale:}
\begin{itemize}[leftmargin=1.4cm, itemsep=1pt]
  \item \textbf{Stop-the-world checkpointing} imposes regular pauses on all processes.
  \item \textbf{Checkpoint overhead} grows with application size; writing 100s of GB is
    slow.
  \item \textbf{Wasted work:} on failure, all processes restart even those that did not fail;
    all computation since the last checkpoint is lost.
  \item \textbf{Scalability crisis:} in a 10,000-node system where each node fails once
    per year, a failure occurs roughly every \textbf{53 minutes} --- far shorter than a
    typical checkpoint interval, making the system permanently stalled.
\end{itemize}
\end{solbox}

%-------
\subsection{Q1(b) --- MapReduce advantage over MPI for failure handling
  \texorpdfstring{\qmarks{2}}{}}
\begin{qbox}What advantage does MapReduce have over MPI-based systems for failure handling and
recovery?\end{qbox}
\begin{solbox}
\begin{enumerate}[leftmargin=1.6cm, itemsep=5pt]
  \item \textbf{Task-level re-execution, not global restart.}\\
    When a MapReduce worker fails, only the \textbf{specific Map or Reduce tasks} running on that
    worker are re-scheduled on another available node. All other in-progress and completed tasks
    continue unaffected. Contrast with MPI: \textbf{all} processes restart from checkpoint.

  \item \textbf{No application-level checkpointing needed.}\\
    Map/Reduce tasks are \textbf{stateless and deterministic}: their output depends only on their
    input (read from the replicated distributed FS). Any failed task can be re-run from the
    original DFS input --- which is still available because GFS/HDFS keeps \textbf{3 replicas}.
    There is no global state to save or restore.

  \item \textbf{Failure handling is transparent to the programmer.}\\
    The framework (master + YARN) detects failures via \textbf{heartbeats} and automatically
    re-queues tasks. The developer writes zero fault-tolerance code.
\end{enumerate}

\textbf{Summary:} MPI requires saving and restoring the \textit{entire global state}; MapReduce
re-runs only the \textit{individual failed tasks} --- a fundamentally more scalable fault model.
\end{solbox}

%-------
\subsection{Q1(c) --- GFS replica location API: MapReduce optimization enabled
  \texorpdfstring{\qmarks{2.5}}{}}
\begin{qbox}GFS exposes the locations of block replicas to clients via an API. What important
MapReduce optimization does this enable?\end{qbox}
\begin{solbox}
The optimization enabled is \textbf{Data Locality} (also called
\textit{moving computation to data} or \textit{locality-aware scheduling}).\\[6pt]

\textbf{How it works:}
\begin{enumerate}[leftmargin=1.6cm, itemsep=3pt]
  \item The MapReduce master \textbf{queries the GFS API} to learn which servers hold each
    input block's replicas.
  \item It then \textbf{schedules each Map task preferentially} on:
    \begin{enumerate}[leftmargin=1.2cm, itemsep=1pt]
      \item[1.] A server that already holds a replica of that block (\textbf{local disk read}), or
      \item[2.] A server on the \textbf{same rack} as one holding a replica (intra-rack read), or
      \item[3.] Any server anywhere (fallback --- full cross-rack network read).
    \end{enumerate}
  \item Most Map tasks read their input data \textbf{from local disk} instead of over the network.
\end{enumerate}

\textbf{Why it matters:}
\begin{itemize}[leftmargin=1.4cm, itemsep=1pt]
  \item Eliminates network bandwidth as a bottleneck for the Map phase (which reads the entire
    input dataset).
  \item In a large cluster, \textbf{network cross-section bandwidth} is the scarcest resource;
    local reads conserve it for the shuffle phase (which must traverse the network).
  \item Measured speedup: significantly faster job completion at scale (often the dominant
    optimisation in MapReduce performance).
\end{itemize}
\end{solbox}

%-------
\subsection{Q1(d) --- Two limitations on the Map function (repeat)
  \texorpdfstring{\qmarks{2}}{}}
\begin{qbox}What are two important limitations that MapReduce places on the Map function?\end{qbox}
\begin{solbox}
\textit{(Same question as 2022 Q4(a) --- see full answer there.)}\\[4pt]

\textbf{Summary (two limitations):}
\begin{enumerate}[leftmargin=1.6cm, itemsep=4pt]
  \item \textbf{No side effects --- must be a pure/deterministic function.}
    Map must produce the same output for the same input; it cannot write to shared external state.
    This is required for safe re-execution on failure and for speculative execution.
  \item \textbf{No inter-task communication --- each record processed independently.}
    A Map function sees only one input record at a time; it cannot access other records or the
    outputs of sibling Map tasks. This is what makes MapReduce trivially parallelisable.
\end{enumerate}
\end{solbox}

%=====================================================================
\section{Quick Summary --- Exam Patterns \& Must-Know}
%=====================================================================
\begin{notesbox}
\begin{itemize}[leftmargin=1.3cm, itemsep=3pt]
  \item \textbf{2024 entire Q1 = MapReduce} (4 parts, 9 marks of Section A).
    Trend: MapReduce yield has jumped from 3/5 to 4/5. Treat it as a primary topic.

  \item \textbf{Two Map limitations (2022 Q4a = 2024 Q1d, verbatim repeat):}
    (1) No side effects / deterministic.
    (2) No inter-task communication / independent records.
    Know both with the WHY (re-execution + parallelism).

  \item \textbf{Straggler / Speculative Execution (2022 Q4b):}
    Define straggler $\rightarrow$ barrier problem $\rightarrow$ track progress \%
    $\rightarrow$ launch backup copy $\rightarrow$ use whichever finishes first.

  \item \textbf{MPI vs.\ MapReduce fault tolerance (2024 Q1a + Q1b):}
    MPI = global checkpoint/restart (all processes stop and roll back).
    MapReduce = re-run only failed task (task-level, stateless, DFS replicas survive node death).

  \item \textbf{GFS locality (2024 Q1c):}
    GFS API reveals replica locations $\rightarrow$ master schedules Map tasks co-located with data
    $\rightarrow$ eliminates network reads in Map phase $\rightarrow$ major performance gain.
    Priority: same node $\to$ same rack $\to$ anywhere.

  \item \textbf{Types of MapReduce jobs (2021 Q7c):}
    Word count, distributed grep, reverse web-link graph, URL frequency, distributed sort,
    inverted index. Know Map/Reduce steps for at least word count and grep cold.

  \item \textbf{YARN (not directly examined yet, but know components):}
    Resource Manager (global scheduler) + Node Manager (per-server daemon)
    + Application Master (per-job; detects task failures).

  \item \textbf{Key pipeline to draw:}
    Input (DFS) $\to$ Map (local disk output) $\to$ Shuffle/hash-partition
    $\to$ Reduce (DFS output). Barrier: all Maps must finish before Reduce starts.
\end{itemize}
\end{notesbox}

\end{document}
